\chapter{Explicit Effects}\label{chap:effects}
% Budget: 11 standard pages. The theory of the solution. Statements and short arguments here, full proofs
% in Appendix~\ref{app:proofs}. Material: the deleted chapters/explicit-effects.tex in git history
% (commit before 24 Oct), research notebook F31--F43.

\section{A Language with Heap Tokens}\label{sec:effects-language}
% Syntax, static single use, token typing, the Chains lemma.

\section{Semantics}\label{sec:effects-semantics}
% Big-step semantics, the trap model, erasure: tokens carry nothing.

\section{Partitions}\label{sec:effects-partitions}
% Field, object and object-field partitions from points-to facts; the refinement order.

\section{Signatures}\label{sec:effects-signatures}
% The least solution over the call graph; the signature theorem; the frame lemma.

\section{What the Tokens Guarantee}\label{sec:effects-guarantees}
% Locality and independence; forwarding and dead stores as consequences.

\section{Exceptions and Deoptimization}\label{sec:effects-exceptions}
% Calls give back an exception and branch; traps only for deoptimization.

\section{Concurrency}\label{sec:effects-concurrency}
% Acquire writes and release reads the shared partitions; thread escape decides what is shared;
% synchronization on objects only one thread reaches orders nothing.

\section{From Signatures to Kill Sets}\label{sec:effects-kill-sets}
% The kill-set theorem under Graal's memory model; the virtual machine's state; why no call kills init.
